% =====================================================================================
% Chapter 6 · Funnel mediation — how a treatment works
%
% Built from notebooks/04_funnel_mediation.ipynb. NOT ONE NUMBER IN THIS FILE IS TYPED BY
% HAND: every figure in the prose is a macro from build/macros.tex, generated by cmp.report
% from the executed notebook. The only literals below are (i) the structural coefficients of
% the data-generating model in §6.2, which are *definitions* rather than results, and (ii) the
% conventions stated in the text (the 90 % interval level, the 0.05 test size of the
% falsification budget), which are choices, not measurements.
% =====================================================================================
\chapter{Funnel mediation: how a treatment works}
\label{ch:med}

\begin{quote}\small
\emph{A growth team has one budget and four funnel stages --- onboarding, engagement,
activation, conversion --- and a conversion is worth \eur{\nbFourValuePerConv}. A better
onboarding flow raises conversion by \nbFourBayesTotal{} conversions per unit of onboarding
score. The budget question is not whether it works but \emph{how}: money spent on onboarding is
worth spending only if the mechanism it fires is the one the firm believes in. This chapter
decomposes that total effect and finds it is essentially all \emph{indirect} --- a natural
indirect effect of \nbFourBkNie{} against a natural direct effect of \nbFourBkNde{}, whose
interval straddles zero, so \nbFourPropMediated\,\% of onboarding's value is realised two stages
downstream and a last-touch attribution model would price the lever at roughly nothing. The
classical arm gets there with two regressions and a multiplication; the bootstrap the applied
literature insists on turns out to agree with the delta method to \nbFourSeGapPct\,\%, and the
posterior --- which needs neither --- returns intervals that are, if anything, \emph{wider}
(width ratios \nbFourWidthMin--\nbFourWidthMax). The uncertainty that matters is not in any of
those intervals. It is in \cref{as:med:si2}, which no experiment on the treatment can buy and no
resampling can see: \cref{sec:med:wrong} shows a hidden mediator--outcome confounder explaining
\nbFourFlipMed\,\% of one fitted edge overturning the investment ranking, while both apparatuses
report their tight intervals with the same untroubled face.}
\end{quote}

% =====================================================================================
\section{The decision}
\label{sec:med:decision}

A subscription business acquires users, walks them through an onboarding flow, hopes they engage
with the product, hopes engagement carries them across an activation threshold, and finally sells
them something. Four stages, one budget. The growth team must decide where the next euro goes,
and the stages are not independent bets: money spent on onboarding does not stay in onboarding.
It cascades. A better first-run experience produces more engaged users, who activate more often,
who convert more often. The euro is spent at the top and paid back at the bottom.

That cascade is what makes the naive answer wrong in a specific and expensive way. Rank the
stages by their correlation with revenue and onboarding will come last, because by the time a user
converts, onboarding's contribution has been laundered through two intermediate behaviours that
correlate with conversion far more tightly than onboarding does. Worse --- and this is the failure
that actually happens in production --- put the funnel stages into one regression as
``controls'' and onboarding's coefficient collapses to \nbFourNaiveOnCoef{} against a true total
effect of \nbFourTrueTotal. The regression has not measured onboarding's value; it has
\emph{conditioned it away}, by holding fixed the very channels through which onboarding acts.
Every last-touch and last-click attribution system in commercial use commits a version of this
error, and it systematically defunds the top of the funnel.

So the question this chapter answers is not ``does onboarding work?'' --- \cref{sec:med:classical}
settles that in one line --- but \textbf{through what?} The total effect of a treatment splits into
the part that travels through a named intermediate variable, the \emph{mediator}, and the part
that does not. That split is the object the budget consumes. If onboarding's value is realised
downstream, it is a \emph{root} investment whose payoff must be credited two stages away, and any
system that scores it on direct conversions will underfund it. If instead a large share of the
effect is direct, onboarding is doing something the funnel model does not capture, and the funnel
model is the thing to fix.

The wrong answer is expensive in both directions and, unusually for this book, the expense is not
symmetric between the two arms of the analysis --- because as \cref{sec:med:compare} shows, the
classical and Bayesian arms give the \emph{same} answer, right or wrong. The risk lives entirely
in \cref{sec:med:ident}. Mediation asks a strictly harder question than the rest of this book, and
buys the answer at a strictly higher price in assumptions. The chapter's purpose is to make that
price visible, and then to pay it in public.

% =====================================================================================
\section{The data-generating model}
\label{sec:med:dgm}

A causal method cannot be graded on real data, because on real data the split is exactly what is
missing. So the method is first graded against a world whose paths are known by construction.

The simulator draws $n = \nbFourNObs$ users. Two variables are \emph{exogenous} --- nothing in the
model points into them: the onboarding score $O$ that a user's flow achieves, and the
\emph{channel quality} $Q$ of the acquisition source that brought them in. Each is uniform on
$[0,1]$. The three downstream stages are structural equations, one per arrow of
\cref{fig:med:dag}:
%
\begin{align}
  O,\ Q &\;\sim\; U(0,1), \label{eq:med:exog}\\
  \text{engagement} &\;=\; 2.2\,O \;+\; 1.1\,Q \;+\; \varepsilon_1,
        &\varepsilon_1 &\sim \mathcal N(0,\ 0.4^{2}), \label{eq:med:eng}\\
  \text{activated} &\;=\; 1.6\,\text{engagement} \;+\; 0.6\,Q \;+\; \varepsilon_2,
        &\varepsilon_2 &\sim \mathcal N(0,\ 0.4^{2}), \label{eq:med:act}\\
  \text{converted} &\;=\; 0.9\,\text{activated} \;+\; 0.35\,\text{engagement} \;+\; \varepsilon_3,
        &\varepsilon_3 &\sim \mathcal N(0,\ 0.3^{2}). \label{eq:med:conv}
\end{align}
%
Three features of this world are deliberate, and each plants a specific lesson.

\textbf{There is no $O \to \text{converted}$ arrow.} \Cref{eq:med:conv} does not contain $O$. The
true natural direct effect is exactly $\nbFourTrueNde$, and every unit of onboarding's value
travels through the funnel. This is not the assumption the model will be given: \cref{sec:med:bayes}
\emph{fits} that edge anyway, so a near-zero direct effect is a discovery rather than an artefact of
omitting the arrow.

\textbf{The paths are products.} Because the system is linear, the effect along a route is the
product of the coefficients on its arrows. Onboarding reaches conversion by two routes: the long
one, $O \to \text{eng} \to \text{act} \to \text{conv}$, worth
$2.2 \times 1.6 \times 0.9 = \nbFourTrueViaAct$; and the short one, $O \to \text{eng} \to
\text{conv}$, worth $2.2 \times 0.35 = \nbFourTrueViaEng$. Their sum is the natural indirect
effect, and since the direct effect is zero, it is also the total: $\nbFourTrueTotal$. The planted
proportion mediated is therefore $100\,\%$. Those are the numbers every estimator in this chapter
is marked against.

\textbf{Channel quality feeds \emph{two} stages.} $Q$ appears in \cref{eq:med:eng} and again in
\cref{eq:med:act}. It is therefore a common cause of two mediators, and if it is not adjusted for,
the engagement-to-activation edge absorbs its association and every path product built on that
edge is poisoned. This is not a decorative complication. It is the single most realistic way a
funnel analysis goes wrong in production --- the growth team never logged acquisition channel ---
and \cref{sec:med:wrong} deletes $Q$ and measures the damage in both arms.

Note what is \emph{absent} from \crefrange{eq:med:eng}{eq:med:conv}: any interaction term
$O \times \text{engagement}$. That absence is what makes the two-way split of
\cref{sec:med:ident} add up exactly, and \cref{as:med:nointer} states it as an assumption rather
than smuggling it in as an algebraic convenience.

\input{build/floats/nb04_dag}

% =====================================================================================
\section{Identification}
\label{sec:med:ident}

\subsection*{The estimand}

Write $X$ for the treatment (onboarding score), $M$ for the mediator, $Y$ for the outcome
(conversion) and $C$ for the observed common causes (here $Q$). In the potential-outcome notation
of \citet{rubin1974}, mediation needs a \emph{two-argument} potential outcome: $Y(x, m)$ is the
conversion a user would show if their onboarding were set to $x$ \emph{and} their mediator to $m$,
and $M(x)$ is the mediator value that onboarding $x$ would produce. The total effect of moving
onboarding from $x'$ to $x$ is then $\E[Y(x, M(x)) - Y(x', M(x'))]$, and \citet{pearl2001} and
\citet{robins1992} split it as
%
\begin{align}
  \text{NDE} &\;=\; \E\big[\,Y\big(x,\,M(x')\big) \;-\; Y\big(x',\,M(x')\big)\,\big],
    \label{eq:med:nde}\\
  \text{NIE} &\;=\; \E\big[\,Y\big(x,\,M(x)\big) \;-\; Y\big(x,\,M(x')\big)\,\big],
    \label{eq:med:nie}\\
  \text{TE}  &\;=\; \text{NDE} \;+\; \text{NIE}. \label{eq:med:te}
\end{align}
%
Read them slowly, because the notation is doing something no experiment ever does. The
\emph{natural direct effect} of \cref{eq:med:nde} moves the treatment while pinning the mediator at
the value it \emph{would have taken under the other treatment}. The \emph{natural indirect effect}
of \cref{eq:med:nie} does the reverse: it fixes the treatment and swaps which world the mediator
comes from. Both are \textbf{cross-world} quantities. $Y(x, M(x'))$ combines a treated outcome with
an untreated mediator, and \emph{no single individual, in any experiment, ever exhibits that
combination}. The two arguments live in different worlds. That is what makes the assumptions below
strong, and it is the reason \cref{as:med:crossworld} can never be tested by anybody, ever.

The derived quantity the business quotes is the \emph{proportion mediated},
$\text{NIE}/\text{TE}$, and this chapter reports it with an interval, because a headline of ``one
hundred per cent mediated'' is a ratio of two estimated quantities and is not exempt from
uncertainty merely because people like to say it out loud.

\subsection*{What buys the decomposition}

\begin{assumption}[Treatment ignorability]\label{as:med:si1}
$\{\,Y(x,m),\; M(x')\,\} \perp X \mid C$: given the observed common causes, who received good
onboarding is as-good-as-random with respect to both the outcome and the mediator. \emph{Violation:}
motivated users breeze through onboarding \emph{and} would have converted anyway, and onboarding
inherits motivation's credit. \emph{Status here:} true by construction --- \cref{eq:med:exog}
draws $O$ independently of everything --- and \emph{it is the only one of the four an experiment
can buy}. \Cref{sec:med:wrong} shows the total effect surviving intact when the analysis is broken
in other ways, precisely because this assumption holds.
\end{assumption}

\begin{assumption}[Mediator ignorability]\label{as:med:si2}
$Y(x,m) \perp M \mid X = x,\ C$: given treatment and the observed common causes, the mediator is
as-good-as-random with respect to the outcome. \emph{Violation:} a user's intrinsic interest in the
product drives both their engagement \emph{and} their conversion; the fitted
engagement-to-conversion edge then contains that trait's work, $\hat b$ is not the causal edge, and
the direct/indirect \emph{split} is wrong even when the total is fine. \emph{Status here:}
\textbf{untestable}. It is the assumption on which this chapter's entire decision rests, and
\cref{sec:med:wrong} does the only honest thing available --- it prices what a violation would cost.
\end{assumption}

\begin{assumption}[Positivity]\label{as:med:pos}
$0 < p(M = m \mid X, C)$ for every mediator value $m$ in the support: every mediator level must be
reachable at every treatment level. \emph{Violation:} if no poorly-onboarded user ever reaches high
engagement, the counterfactual ``high engagement without onboarding'' of \cref{eq:med:nie} is pure
extrapolation from the model's functional form. \emph{Status here:} satisfied by construction ---
the Gaussian noise of \cref{eq:med:eng} gives every mediator value positive density at every $O$ ---
and on real data it is checkable, by plotting the mediator's conditional distribution across the
treatment range.
\end{assumption}

\begin{assumption}[Cross-world independence]\label{as:med:crossworld}
$Y(x,m) \perp M(x') \mid C$: there is no mediator--outcome confounder that is \emph{itself caused by
the treatment}. \emph{Violation:} onboarding creates a habit which separately drives both engagement
and conversion. \emph{Status here:} \textbf{untestable even in principle}, and not merely in
practice: the two potential outcomes it constrains belong to different worlds and cannot be jointly
observed in any experiment, however large or however well randomized \citep{robins1992, imai2010}.
It is the honest price of asking a cross-world question, and the reason some authors prefer
\emph{controlled} direct effects, which require no such condition, at the cost of no longer
decomposing the total.
\end{assumption}

\begin{assumption}[No mediator--mediator confounding]\label{as:med:mm}
With two mediators in sequence, ``the effect through activation'' is a \emph{path-specific} effect
\citep{avin2005}, and identifying it requires in addition that the two mediators share no unmeasured
common cause \citep{vanderweele2014mediators}. \emph{Violation:} in this world, $Q$ --- which drives
both of them. \emph{Status here:} \textbf{partially testable}. The graph
implies conditional independences, and \cref{sec:med:valid} checks them: of \nbFourImpTests{}
implied independences, \nbFourImpViolations{} is rejected at the $5\,\%$ level, which a Binomial
false-positive budget prices as unremarkable ($\Prob \ge \nbFourImpViolations$ by chance
$= \nbFourImpPchance$). And when the assumption is deliberately broken in \cref{sec:med:wrong}, the
same test fires: \nbFourBadImpViolations{} of \nbFourBadImpTests. This is the one strong assumption
of the chapter that leaves a testable footprint, and the falsification cell is not decoration.
\end{assumption}

\begin{assumption}[No treatment--mediator interaction]\label{as:med:nointer}
The outcome equation contains no $X \times M$ term, so the return to engagement does not depend on
how well the user was onboarded. \emph{Consequence:} the additive decomposition of \cref{eq:med:te}
holds exactly, and the natural effects do not depend on the baseline treatment level $x'$.
\emph{Status here:} true by construction (\cref{eq:med:conv}), and \textbf{not tested} --- an
honesty this chapter would rather state than bury. On real data it is plausible that better-onboarded
users engage more \emph{productively}, in which case NDE $+$ NIE $\neq$ TE and the two-way split must
give way to \citeauthor{vanderweele2014}'s four-way decomposition into direct, mediated,
interactive and mediated-interactive components \citep{vanderweele2014, vanderweele2015}. The
correct move is to fit the interaction and test it, not to assume it away.
\end{assumption}

\begin{remark}[Why an RCT does not rescue you]
This is the sentence the chapter exists to deliver. Randomizing the treatment buys
\cref{as:med:si1} and \emph{nothing else}. It does not touch \cref{as:med:si2}, because
\textbf{randomizing onboarding does not randomize engagement}: engagement is a behaviour the user
chooses, not a knob the experimenter turns. In a perfect randomized trial of an onboarding flow,
with a hundred thousand users and flawless compliance, the total effect is identified by the coin
flip and the \emph{split} is identified by nothing at all \citep{bullock2010, imai2011}. A
mediation analysis inside an RCT is an observational study wearing an experiment's clothes, and
\crefrange{as:med:si2}{as:med:crossworld} are as unearned there as they are here. What an experiment
\emph{would} buy is a design that randomizes the mediator too --- a second, cross-randomized arm, or
an instrument for engagement --- and that is a different and more expensive study than the one the
growth team ran.
\end{remark}

\subsection*{The mediation formula, and why it becomes a product}

Under \crefrange{as:med:si1}{as:med:crossworld}, the cross-world mean of
\crefrange{eq:med:nde}{eq:med:nie} --- a quantity nobody observes --- is identified by observed
conditional distributions. This is \citeauthor{pearl2001}'s \emph{mediation formula}
\citep{pearl2001, imai2010}:
%
\begin{equation}
  \E\big[\,Y\big(x,\,M(x')\big)\,\big] \;=\;
  \int_c \int_m \E\big[\,Y \mid X = x,\ M = m,\ C = c\,\big]\;
  \mathrm dP\big(m \mid X = x',\ C = c\big)\; \mathrm dP(c).
  \label{eq:med:formula}
\end{equation}
%
Read \cref{eq:med:formula} as a recipe: predict the outcome under treatment $x$, but average over the
mediator distribution that treatment $x'$ \emph{would have produced}. The inner integral is what
crosses the worlds, and the assumptions are what license it. It is estimator-agnostic --- nothing in
it commits to linearity or to any particular fitting procedure.

Specialise it to a linear system. Let $M = \alpha_0 + a X + \gamma C + \varepsilon_M$ and
$Y = \beta_0 + c X + b M + \delta C + \varepsilon_Y$, with mean-zero noise and no interaction
(\cref{as:med:nointer}). Then $\E[Y \mid X = x, M = m, C] = \beta_0 + c\,x + b\,m + \delta C$ and
$\E[M \mid X = x', C] = \alpha_0 + a\,x' + \gamma C$, so \cref{eq:med:formula} collapses to
$\E[Y(x, M(x'))] = \text{const} + c\,x + b\,a\,x'$, and hence, for a unit contrast
$\Delta x = 1$,
%
\begin{equation}
  \boxed{\ \text{NDE} \;=\; c, \qquad \text{NIE} \;=\; a\,b, \qquad \text{TE} \;=\; c + a\,b.\ }
  \label{eq:med:product}
\end{equation}
%
That is the classic ``$a \cdot b$'' of the mediation literature \citep{baron1986}, now
\emph{derived} from \cref{eq:med:formula} rather than asserted. The distinction matters: the product
of coefficients is not a rival theory of mediation but the mediation formula's linear special case,
which is precisely why \cref{sec:med:classical} can reach the chapter's answer with two regressions
and no likelihood at all.

\begin{remark}[Why \citeauthor{baron1986} is not enough]
\Citeauthor{baron1986}'s 1986 paper is the most cited article in the history of social psychology
and it gets the algebra right, but it predates the potential-outcome machinery above and therefore
cannot say what its own product \emph{estimates}. Three consequences. First, it offers no estimand:
$a b$ is a pair of regression slopes multiplied together, and it becomes the NIE of
\cref{eq:med:nie} only under \crefrange{as:med:si1}{as:med:crossworld}, which the paper does not
state. Second, its ``causal steps'' procedure --- require a significant total effect, then a
significant $a$, then a significant $b$ --- is a poor test, with lower power than testing $ab$
directly \citep{mackinnon2002}, and it wrongly rules out mediation whenever a direct and an indirect
path cancel. Third, it says nothing about how to build an interval around a \emph{product}, which is
\cref{sec:med:classical}'s next problem. \Citet{imai2010} and \citet{vanderweele2015} supply all
three missing pieces, and this chapter uses \citeauthor{baron1986}'s arithmetic under
\citeauthor{imai2010}'s licence.
\end{remark}

% =====================================================================================
\section{The classical baseline}
\label{sec:med:classical}

Before any sampler, the discipline of this book is to ask what a competent analyst produces
\emph{without} one. Here the answer is not a different analysis. It is the \emph{same} estimands of
\crefrange{eq:med:nde}{eq:med:te}, bought by the \emph{same} assumptions of \cref{sec:med:ident},
estimated with the oldest tool in the literature: two ordinary least-squares regressions and a
multiplication.

\subsection*{Two regressions and a multiplication}

With $X$ = onboarding score, $M$ = engagement, $Y$ = conversion and $C$ = channel quality, fit
%
\begin{align}
  M_i &= \alpha_0 + a\,X_i + \gamma\,C_i + \varepsilon^M_i, \label{eq:med:ols1}\\
  Y_i &= \beta_0 + c\,X_i + b\,M_i + \delta\,C_i + \varepsilon^Y_i, \label{eq:med:ols2}
\end{align}
%
and read \cref{eq:med:product} straight off the coefficients. The fit gives
$\hat a = \nbFourBkA$ (SE $\nbFourBkASe$, $t = \nbFourBkAT$) and $\hat b = \nbFourBkB$
(SE $\nbFourBkBSe$, $t = \nbFourBkBT$). Both links are strong, which will matter in a moment.
The decomposition is in \cref{tab:med:classical}: a natural indirect effect of \nbFourBkNie{}
against the planted \nbFourTrueTotal, a natural direct effect of \nbFourBkNde{} whose interval
$[\nbFourBkNdeLo, \nbFourBkNdeHi]$ contains zero, and a proportion mediated of
\pc{\nbFourBkProp}. Every interval covers its planted value. \textbf{This is a complete and correct
answer to the chapter's question, and it took milliseconds.}

\input{build/tables/nb04_classical}

\subsection*{The regression you must not run}

One choice in \cref{eq:med:ols2} deserves its own paragraph, because getting it wrong is the most
common way an applied mediation analysis is quietly wrong. The funnel has \emph{two} ordered
mediators, and the single-mediator read takes $M$ = engagement and deliberately leaves
\code{activated} \textbf{out} of the outcome regression --- even though every analyst's instinct
says activation is an obvious control for conversion.

It is not a control. It sits \emph{on the path}
$\text{engagement} \to \text{activation} \to \text{conversion}$, and conditioning on it
\textbf{blocks that path}. Then $\hat b$ stops being engagement's total effect on conversion --- the
reduced form $0.9 \times 1.6 + 0.35 = \nbFourBReduced$ in this world --- and collapses to its
direct-only slice, the planted $\nbFourBDirect$. The numbers oblige: with activation left out,
$\hat b = \nbFourBlockingBOk$; with activation controlled, $\hat b = \nbFourBlockingBBad$, a factor
of \nbFourBlockingFactor. The indirect effect would then be reported as \nbFourBlockingNie{}
instead of \nbFourBkNie: a funnel investment priced at a fifth of its worth.

The rule generalises, and it is the mirror image of \cref{sec:med:decision}'s naive-regression trap:
\textbf{never condition on a variable that lies on the path you are trying to measure.} There, the
error was sprung on the treatment; here, on the mediator. Both are the same mistake, and both are
committed by an analyst doing what feels like due diligence. \Cref{sec:med:wrong} adds the third
member of the family --- the collider --- because conditioning on a post-treatment variable is never
free: it can \emph{open} a path as easily as close one.

\subsection*{Where the interval comes from: the product problem}

$\hat a \hat b$ is a \emph{product} of two estimates. Even when $\hat a$ and $\hat b$ are each exactly
normal, their product is not: it is skewed, and its variance has no exact elementary form. Two
classical devices exist, and the chapter runs both.

The first is the \textbf{delta method}, which linearises the product around the estimates and reads
off the variance of the linearisation \citep{sobel1982}:
%
\begin{equation}
  \text{SE}_\Delta(\hat a \hat b) \;\approx\;
  \sqrt{\ \hat b^{2}\,\text{SE}(\hat a)^{2} \;+\; \hat a^{2}\,\text{SE}(\hat b)^{2}\ }
  \;=\; \nbFourSeDelta .
  \label{eq:med:delta}
\end{equation}
%
It is a \emph{first-order} approximation --- \cref{eq:med:delta} is a one-term Taylor expansion, and
it drops the $\text{Var}(\hat a)\text{Var}(\hat b)$ term entirely --- and it asserts normality for a
quantity that is not normal.

The second is the \textbf{nonparametric bootstrap} \citep{efron1979, efron1993}: resample the
\emph{rows} of the dataset with replacement, refit \emph{both} of
\crefrange{eq:med:ols1}{eq:med:ols2} on each replicate, recompute $\hat a\hat b$, and take
percentiles of the resulting \nbFourNBoot{} values. It assumes only that the rows are exchangeable.
This is what the applied literature has recommended for twenty years \citep{preacher2008,
mackinnon2002}, and it is what \cref{tab:med:classical}'s intervals use.

\input{build/tables/nb04_scales}

\subsection*{The bootstrap was worth running, and it barely mattered}

Both halves of that sentence are true, and the second is the interesting one. \Cref{tab:med:scales}
sets the two devices against each other: the delta method says \nbFourSeDelta, the bootstrap says
\nbFourSeBoot{} --- a ratio of \nbFourSeRatio, agreeing to \pc{\nbFourSeGapPct}.
\Cref{fig:med:product} shows why. The bootstrap sampling distribution of $\hat a\hat b$ has a sample
skew of only \nbFourNieBootSkew. It is, to the eye and to the third decimal, normal.

\input{build/floats/nb04_product}

This is not a failure of the chapter's argument; it is the argument, stated precisely. \textbf{The
delta method does not fail here, and the chapter will not pretend it does.} It fails --- and the
product's skew bites --- when one of the two links is \emph{weak}. A small $\hat a$ drags the
product's distribution off normality (in the limit $a = 0$, $\hat a \hat b$ is a product of two
mean-zero normals, whose density is sharply peaked and heavy-tailed, and nothing like the Gaussian
of \cref{eq:med:delta}), and the delta interval then covers at less than its nominal rate
\citep{mackinnon2002, loeys2015}. Here, with $n = \nbFourNObs$ and $t$-statistics of \nbFourBkAT{}
and \nbFourBkBT{} on the two links, the estimator sits about as far from that regime as it is
possible to sit.

The bootstrap's value in this run is therefore not a \emph{correction}. It is the fact that
\textbf{you could not have known any of the above in advance}: the diagnosis --- both links far from
zero, $n$ large, skew negligible --- is available only once the bootstrap has been run. And the
regime in which the delta method breaks (a marginal first link) is exactly the regime in which the
analyst most wants a trustworthy interval, because a marginal first link is what an inconclusive
mediation story looks like. Bootstrap by default; be pleased, not disappointed, when it agrees.

\subsection*{The whole decomposition, chained}

The single-mediator read collapses the funnel into one hop. To recover the \emph{path-specific}
effects --- the long route through activation and the short route that skips it --- the same
arithmetic is chained over three stage regressions, one per arrow of
\crefrange{eq:med:eng}{eq:med:conv}, and multiplied along each route. Here activation \emph{does}
belong on the right-hand side of the conversion equation, because the goal has changed: the analysis
now \emph{wants} to split the long route from the short one rather than to measure engagement's
total. The results are the classical columns of \cref{tab:med:compare}, and
\nbFourClNCover{} of \nbFourNEstimands{} bootstrap intervals cover their planted value.

\subsection*{What this apparatus cannot say}

Two things, and the second is the one that decides the chapter.

It cannot produce a probability \emph{about the effect}. \Cref{sec:med:euro} must answer ``what is
the probability engagement's euro leverage beats activation's?'' and ``what is the probability this
lever's return on investment clears 1?''. Those sentences ask for probabilities about the parameters
themselves; a confidence interval is a property of a procedure and attaches none. That gap is
familiar from every chapter of this book, and \cref{sec:med:bayes} exists to fill it.

And it cannot see \cref{as:med:si2}. This one is not a gap in the frequentist toolkit --- it is a
gap in \emph{both} toolkits, and it will not be filled by anything in \cref{sec:med:bayes} either.
Resampling the rows re-estimates the same quantity a thousand times over and reports, very
precisely, how stable it is. If a hidden trait drives both engagement and conversion, $\hat b$ is not
the causal edge, $\hat a \hat b$ is not the NIE, and the bootstrap will wrap a tight, confident,
wrong interval around it. Hold that thought until \cref{sec:med:compare}, which measures the
consequence, and \cref{sec:med:wrong}, which prices it.

% =====================================================================================
\section{The Bayesian model}
\label{sec:med:bayes}

The chain of \crefrange{eq:med:eng}{eq:med:conv} is handed to the sampler \emph{whole} --- one
Gaussian likelihood per stage, fitted jointly --- rather than as three disconnected regressions.
Each equation carries an intercept the generating model does not have (harmless: its posterior sits
at zero), and, crucially, the conversion equation carries a \emph{direct} onboarding edge $c_{on}$
which \cref{eq:med:conv} does not contain. It is fitted so that ``how much is mediated?'' can be
answered by the data rather than by the modeller's choice of which arrows to draw:
%
\begin{align}
  \text{eng}_i &\sim \mathcal N\big(\beta_0^{(1)} + e_{on}O_i + e_{ch}Q_i,\ \sigma_1^2\big),
    \label{eq:med:lik1}\\
  \text{act}_i &\sim \mathcal N\big(\beta_0^{(2)} + a_{eng}\,\text{eng}_i + a_{ch}Q_i,\
    \sigma_2^2\big), \label{eq:med:lik2}\\
  \text{conv}_i &\sim \mathcal N\big(\beta_0^{(3)} + c_{act}\,\text{act}_i
    + c_{eng}\,\text{eng}_i + c_{on}O_i,\ \sigma_3^2\big), \label{eq:med:lik3}\\
  \beta_\bullet &\sim \mathcal N(0,\ 10^2), \qquad
  \sigma_\bullet \sim \text{HalfNormal}(1). \label{eq:med:prior}
\end{align}
%
The priors of \cref{eq:med:prior} are weakly informative and tied to the scale of the data: $O$ and
$Q$ live on $[0,1]$, the stages are small multiples of them, the planted coefficients are all below
$2.5$ and the noise scales are $0.3$--$0.4$. So $\mathcal N(0, 10^2)$ is diffuse relative to any
plausible effect, and $\text{HalfNormal}(1)$ is generous for the noise. As throughout this book, the
prior is doing almost nothing and the likelihood is doing everything --- and ``almost nothing'' is
checkable rather than assertable: doubling every prior scale to $\mathcal N(0, 20^2)$ and refitting
moves the largest posterior mean by \nbFourPriorShift{} --- \nbFourPriorShiftSd{} posterior standard
deviations. The data set these numbers.

Inference is by the No-U-Turn sampler \citep{hoffman2014} in PyMC \citep{salvatier2016}:
\nbFourNDraws{} post-warmup draws, $\hat R = \nbFourRhat$, minimum bulk effective sample size
\nbFourEss{} \citep{vehtari2021}, and \nbFourDivergences{} divergent transitions. Nothing that
follows is a sampler artefact.

\subsection*{The one thing the posterior gives away for free}

Here is the structural point of this section, and it is worth stating before any number. The natural
indirect effect of \cref{eq:med:product} is a \emph{deterministic function of two sampled
coefficients}. Given posterior draws $\{a^{(s)}, b^{(s)}, c^{(s)}\}_{s=1}^{S}$ from the joint
posterior, the posterior of the NIE is obtained by
%
\begin{equation}
  \text{NIE}^{(s)} \;=\; a^{(s)}\, b^{(s)},
  \qquad
  \text{proportion mediated}^{(s)} \;=\; \frac{a^{(s)} b^{(s)}}{a^{(s)} b^{(s)} + c^{(s)}} ,
  \label{eq:med:derived}
\end{equation}
%
which is to say: multiply the draws, and divide the draws. \textbf{No delta method. No bootstrap. No
device of any kind.} The uncertainty in the product \emph{is} the uncertainty in the factors,
propagated exactly --- including their correlation, including when $\hat a$ is small, and including
whatever skew the product actually has. The same one-line move produces the proportion mediated, a
\emph{ratio} of two correlated quantities from the same fit, which classically would require its own
bespoke bootstrap whose interval need not cohere with the interval on the numerator.

The reader should be clear about the size of that claim, because \cref{sec:med:compare} will cut it
down. This is \textbf{convenience, not capability}. \Cref{eq:med:derived} does not compute anything
the bootstrap cannot; it computes it without a device, and it keeps every derived quantity coherent
with every other. That is worth having. It is not worth a paragraph of triumphalism, and the
measurements in the next two sections do not support one.

The chain's path effects follow by multiplying along each route --- the long path's posterior is that
of $e_{on} \cdot a_{eng} \cdot c_{act}$, the short path's that of $e_{on} \cdot c_{eng}$ --- and the
direct effect is simply the posterior of $c_{on}$. It comes back at \nbFourBayesDirect, with a
90\,\% credible interval $[\nbFourBayesDirectLo, \nbFourBayesDirectHi]$ that contains zero. The
proportion mediated is \pc{\nbFourPropMediated}, 90\,\% interval
$[\pc{\nbFourPropLo},\ \pc{\nbFourPropHi}]$ (\cref{fig:med:recovery}). The model was free to find a
direct onboarding effect and reports that there is none --- which, since \cref{eq:med:conv} planted
none, is the correct answer, arrived at without being told.

\input{build/floats/nb04_recovery}

% =====================================================================================
\section{Diagnostics and validation}
\label{sec:med:valid}

\subsection*{Arrow by arrow, before path by path}

A path product is only as good as the arrows it multiplies, so the coefficients are graded
individually before anything is multiplied. Every structural coefficient's posterior mean sits within
two posterior standard deviations of its planted value; the worst is
$|z| = \nbFourWorstZ$. The pieces are sound.

\subsection*{One posterior predictive check per equation}

Convergence says the sampler explored the posterior; it says nothing about whether the model can
\emph{generate the data}. And a mediation model is three likelihoods chained together, so its failure
mode is multiplicative: an over-dispersed engagement equation contaminates \emph{every} path product
that runs through engagement, which here is all of them. The check is therefore run once per stage
equation (\cref{fig:med:ppc}). The share of observed values falling inside the per-user 90\,\%
replicate band is \pc{\nbFourPpcEngagement}, \pc{\nbFourPpcActivated} and \pc{\nbFourPpcConverted}
--- calibrated, with no gross misfit anywhere in the chain.

\input{build/floats/nb04_ppc}

\subsection*{Falsifying the graph}

\Cref{as:med:mm} is the chapter's one strong assumption with a testable footprint. The funnel graph
implies a set of conditional independences, and each can be checked as a partial correlation.
\nbFourImpTests{} implied independences are tested at the $5\,\%$ level and \nbFourImpViolations{}
is rejected. That is not evidence against the graph: at size $0.05$, each test rejects $5\,\%$ of the
time even when the graph is \emph{correct}, so the count of violations should be read against a
Binomial false-positive budget, which prices $\Prob(\ge \nbFourImpViolations$ rejections by chance$)$
at \nbFourImpPchance. A single marginal rejection among several tests is expected noise. Treating any
one violation as fatal is how a correct graph gets thrown away.

The test earns its keep in \cref{sec:med:wrong}, where the graph is deliberately broken and it fires:
\nbFourBadImpViolations{} of \nbFourBadImpTests.

\subsection*{Not a lucky seed}

One dataset is one draw of the world, and a method that recovers the truth once may simply have had
friendly noise. So the whole procedure --- simulate, fit, decompose --- is re-run on
\nbFourRepSeeds{} fresh worlds from \crefrange{eq:med:eng}{eq:med:conv}, and the 90\,\% intervals are
graded on how often they cover (\cref{tab:med:replication}). Of \nbFourRepIntervals{} intervals,
\nbFourRepCovered{} cover --- \pc{\nbFourRepCovPct} --- with a largest bias of \nbFourRepMaxBias.
The intervals of \cref{sec:med:bayes} are calibrated, not lucky.

\input{build/tables/nb04_replication}

One structural caveat on reading that table. The four path quantities are read off the \emph{same}
fit, so they do not fail independently: an unlucky dataset tends to shift them together, and misses
therefore cluster on a world rather than scattering one at a time across the columns. A coverage
count of \nbFourRepCovered{} out of \nbFourRepIntervals{} carries less independent information than
its denominator suggests --- which is an argument for more worlds, not for a different statistic.

% =====================================================================================
\section{Point estimate versus posterior}
\label{sec:med:compare}

\input{build/tables/nb04_compare}

\subsection*{On the number, they agree --- and on the interval, they agree too}

\Cref{tab:med:compare} puts both arms on the five estimands of \cref{sec:med:ident}. The largest gap
between a classical plug-in and its posterior mean is \nbFourMaxGap, on effects of size
$\approx \nbFourTrueTotal$. That is not approximate agreement; it is agreement to the third decimal.
This is the Bernstein--von Mises phenomenon in its plainest dress \citep{gelman2013}: with
$n = \nbFourNObs$, a linear Gaussian likelihood and priors of scale $10$ on coefficients of order
$1$, the posterior \emph{is} the likelihood and ordinary least squares sits at its mode. The two arms
are the same estimator wearing different clothes.

Now the widths, which is where a reader who has been sold Bayesian inference expects the payoff. The
credible intervals are \textbf{not narrower}. They are, if anything, slightly \emph{wider}: the ratio
of credible-interval width to bootstrap-interval width runs from \nbFourWidthMin{} to
\nbFourWidthMax{} across the four path estimands. And the same holds for the standard error of the
indirect effect, which \cref{tab:med:scales} priced three ways: the delta method's \nbFourSeDelta,
the bootstrap's \nbFourSeBoot, and the posterior's \nbFourBayesNieSd{} --- three devices, one number.

So the ledger opens with a debit. Say it plainly, because the notebook does and the book will not
soften it: \textbf{the Bayesian machinery did not move the number, and it did not sharpen the
interval.} A reader who came to this chapter expecting mediation to be the place where Bayes rescues
the estimate should take \cref{tab:med:compare} as the correction.

\subsection*{What Bayes did not buy: \texorpdfstring{\cref{as:med:si2}}{Assumption 6.2}}

This is the important one, and it is measured rather than argued.

\Cref{sec:med:wrong} deletes channel quality --- the common cause of \emph{both} mediators
(\cref{as:med:mm}) --- and refits. The Bayesian split breaks: the engagement-to-activation edge
absorbs the missing variable's work and inflates to \nbFourBadAEng{} against a planted $1.60$, and
\nbFourBadNMiss{} of \nbFourNEstimands{} credible intervals now miss the truth. That much is
expected. What \cref{tab:med:failfix} adds is the sentence the chapter owes the reader:
\textbf{the classical arm breaks in the same direction, by the same amount}. The two arms' biased
answers differ by at most \nbFourBadArmGap{} --- three decimal places of agreement in being wrong.

Of course they do. Both fit the same edges under the same assumption. An unmeasured
mediator--outcome confounder is not a statement about sampling variability, and \emph{nothing in
either apparatus is looking for it}. A prior cannot see it. A bootstrap cannot see it: resampling the
rows re-estimates the same biased quantity \nbFourNBoot{} times and reports, with admirable
precision, how stable the bias is. A tight interval around a confounded number is a confident wrong
answer, and both apparatuses will hand you one with the same untroubled face.

That is the deepest lesson of this chapter, and it is a lesson about mediation rather than about
paradigms: \textbf{in a mediation analysis, the dominant uncertainty is not in the interval at all.}
It is in \cref{as:med:si2}, and it is invisible to every statistic either arm computes.

\subsection*{What the posterior did buy}

Four things, and they are real, and none of them is the estimate.

\emph{The product carries its uncertainty without a device.} \Cref{eq:med:derived} is one line, it
needs no linearisation and no \nbFourNBoot{} refits, and --- the part that matters --- it stays
correct in the regime where \cref{eq:med:delta} does not. The delta method's failure mode is a weak
first link; the posterior does not have a failure mode there, because it never approximated anything.
This run happens to be the easy case, and \cref{fig:med:product} shows all three devices agreeing.
The posterior is the arm that did not have to check.

\emph{Derived quantities are functions of draws.} The proportion mediated is a ratio of two
correlated estimates from one fit; \cref{sec:med:euro}'s break-even cost ratio is a ratio of two
lever values; the sensitivity sweep of \cref{sec:med:wrong} is a per-draw transformation of two
edges. Each is one line of arithmetic on the draws, and their intervals cohere with one another by
construction. Classically each is a separate bootstrap, and separate bootstraps do not have to agree
about anything.

\emph{The sensitivity analysis becomes a posterior rather than a point.} \Cref{sec:med:wrong} does
not ask ``what if the engagement-to-conversion edge were partly confounded?'' once, at the point
estimate. It asks it \emph{per draw}, so the flip point comes back as a distribution with an
interval, and the two affected levers shrink \emph{together} because they ride the same edge.

\emph{The decision quantities exist at all.} $\Prob(\text{indirect} > \text{direct}) =
\nbFourPIndGtDir$. $\Prob(\text{proportion mediated} > 90\,\%) = \nbFourPPropGtNinety$.
$\Prob(\text{direct effect} > 0.25) = \nbFourPDirectGtQuarter$. The classical column for each of these
reads \emph{not defined} --- not ``hard'', not ``small'', but undefined, because a confidence
interval attaches no probability to a hypothesis about a fixed parameter. And note the trap that
makes this sharper than it first looks: one \emph{can} bootstrap the proportion mediated, and one
does get an interval --- but the share of bootstrap replicates exceeding $90\,\%$ is \textbf{not}
$\Prob(\text{proportion mediated} > 90\,\%)$. It is the frequency with which a \emph{resampled
estimate} exceeds $90\,\%$: a statement about the estimator's wobble, not about the world.
\Cref{sec:med:euro}'s budget rule needs the latter.

\begin{remark}[The ledger, in one line]
Bayes bought the euro decision of \cref{sec:med:euro} --- the probabilities, the coherent ratios,
the sweep with bands --- and bought nothing at all on the estimate, on the interval, or on the
assumption. In a chapter about mediation, that last clause is the one to remember: the hard part was
never the arithmetic, and no amount of machinery on either side of the aisle will buy you
\cref{as:med:si2}. Only a better measurement will.
\end{remark}

% =====================================================================================
\section{The decision, in euros}
\label{sec:med:euro}

A conversion is worth \eur{\nbFourValuePerConv}. The growth team must now rank three levers ---
onboarding, engagement, activation --- and split a budget across them. Three moves, each correcting
a trap.

\subsection*{Move 1: euros per unit is not a rankable quantity}

The obvious metric is the euro value of a one-unit move in each stage. It is a trap, and the trap is
arithmetic rather than causal. A unit of onboarding score is an enormous move --- the whole variable
has a standard deviation of \nbFourSdOnboarding{} --- while a unit of activation is a fraction of its
standard deviation of \nbFourSdActivated. The two ``units'' are not comparable moves, so a
euros-per-unit ranking is not a ranking of anything; it \emph{reverses} the moment the analyst
switches to comparable per-standard-deviation moves. Both columns are printed side by side in
\cref{tab:med:levers} precisely so the reversal is visible rather than asserted.

The decision-relevant metric is a return on investment: the euro value of a 1-SD lift, divided by the
euro cost of producing a 1-SD lift. At the stated cost of \eur{\nbFourCostSd} per standard deviation,
the ranking is \nbFourRanking, every lever clears break-even
($\Prob(\text{ROI} > 1) = \nbFourProiEngagement$ for the leader), and
$\Prob(\text{engagement beats activation per SD}) = \nbFourPEngBeatsAct$ (\cref{fig:med:forest}).

\input{build/tables/nb04_levers}
\input{build/floats/nb04_forest}

That saturated probability is itself the finding, and it is a warning rather than a
reassurance. With $n = \nbFourNObs$ and a linear model, \emph{sampling} uncertainty is simply not the
binding risk in this decision. The forest plot's intervals are narrow, and they are narrow about the
wrong thing. What the plot cannot show is the possibility that the engagement-to-conversion edge is
partly a hidden confounder's work --- and that is the risk that actually decides the ranking.

\subsection*{Move 2: turn the stated cost into a break-even rule}

``\eur{\nbFourCostSd} per SD for every stage'' is an illustration, not a fact. Real per-stage costs
are precisely what a budget split hinges on, and the analyst does not know them --- finance does. So
rather than bet the recommendation on one invented cost number, the chapter hands finance a rule.
Lever $A$ beats lever $B$ on return exactly when
%
\begin{equation}
  \frac{\text{cost}_A}{\text{cost}_B} \;<\; \frac{\text{value}_A}{\text{value}_B},
  \label{eq:med:breakeven}
\end{equation}
%
and the right-hand side of \cref{eq:med:breakeven} is a \emph{posterior} the analysis already has ---
a ratio of two correlated quantities from the same fit, which is exactly the kind of derived object
\cref{eq:med:derived} produces for free. Onboarding therefore beats engagement on ROI as soon as a
1-SD onboarding lift costs less than \nbFourBeOnbEng$\times$ a 1-SD engagement lift (90\,\%
$[\nbFourBeOnbEngLo,\ \nbFourBeOnbEngHi]$ --- the entire band sits below one). Activation beats
engagement below \nbFourBeActEng$\times$ (90\,\% $[\nbFourBeActEngLo,\ \nbFourBeActEngHi]$).

This is the deliverable, and it is worth more than a verdict. Finance measures the cost ratio, which
is a thing finance can actually measure; the posterior says who wins, with uncertainty attached. The
budget question has been reduced to a measurement the firm is competent to make.

\subsection*{Move 3: the number the CMO will quote, and what it means}

\pc{\nbFourPropMediated} of onboarding's conversion value is mediated
(90\,\% $[\pc{\nbFourPropLo},\ \pc{\nbFourPropHi}]$). Onboarding is a \emph{root} investment whose
payoff lands two stages downstream. Judging it by direct or last-touch conversion --- as the naive
controlled regression of \cref{sec:med:decision} effectively does, returning \nbFourNaiveOnCoef{}
against a true \nbFourTrueTotal{} --- would price it at roughly zero and defund it. That is the
sentence to put in the memo, and it is the chapter's commercial payload.

% =====================================================================================
\section{What can go wrong}
\label{sec:med:wrong}

\subsection*{The unmeasured confounder of the two mediators}

\Cref{as:med:mm} is broken deliberately, in the most realistic way available: the model is refit with
every channel-quality term deleted, exactly as if the growth team had never logged acquisition
channel. \Cref{tab:med:failfix} records the damage in both arms.

The mechanism is predictable in advance and the numbers confirm it. With $Q$ gone, the
engagement-to-activation edge is the only route left through which the model can express the fact
that users from good channels are both more engaged \emph{and} more activated, so it absorbs $Q$'s
work: $a_{eng} = \nbFourBadAEng$ against the planted $1.60$. Every path product built on that edge
inherits the bias --- the long route is overstated by \pc{\nbFourBadViaActErrPct} --- and
\nbFourBadNMiss{} of \nbFourNEstimands{} intervals miss.

Two things are worth extracting from this.

First, \textbf{the raw total effect was never in danger.} A plain regression of conversion on
onboarding alone still returns \nbFourBadOlsTotal{} (SE \nbFourBadOlsSe) against a true
\nbFourTrueTotal, \nbFourBadOlsZ{} standard errors away. Onboarding is exogenous, so omitting $Q$
does not confound the treatment; it confounds the \emph{mediators}. The split breaks and the total
survives, which is precisely the sense in which mediation is harder than the treatment-effect
problems of the surrounding chapters: there, an experiment could in principle buy the assumption;
here, the mediator is a behaviour no experiment assigns.

Second, \textbf{this particular violation left a footprint}, and the graph's own falsification test
found it: \nbFourBadImpViolations{} of \nbFourBadImpTests{} implied independences now fail. That is
genuinely reassuring --- and it must not be over-read. The footprint exists because the omitted
confounder was a common cause of two \emph{observed} mediators, so its absence made the observed data
inconsistent with the graph. A confounder of the mediator and the \emph{outcome} --- \cref{as:med:si2}
--- leaves no such trace, because there is no third observed variable for it to contradict. Which is
the subject of the next subsection.

\input{build/tables/nb04_failfix}

\subsection*{The one that leaves no footprint: pricing \texorpdfstring{\cref{as:med:si2}}{Assumption 6.2}}

An assumption with no accompanying test is a wish. \Cref{as:med:si2} admits no test, so what it gets
instead is a \emph{price}: how large would a hidden mediator--outcome confounder have to be before it
overturned the decision \citep{imai2010, cinelli2020}?

Suppose a hidden trait --- intrinsic product interest, cohort quality --- explains a fraction of the
fitted engagement-to-conversion edge, so that the causally correct edge is smaller than the fitted
one by that share. Two constraints make the sweep honest, and both are easy to get wrong.

It must be \textbf{coherent}. That same edge carries the engagement lever \emph{and} part of
onboarding's total effect. An analysis that shrinks the engagement lever while holding onboarding's
total fixed is not describing any possible world. Both must move together.

It must be \textbf{per draw}. Applying the confounding correction to the point estimate yields a
single flip point, and a single flip point invites the reader to believe it. Applying it inside each
posterior draw yields a \emph{posterior} for the flip point, which is what an honest memo quotes.

\Cref{fig:med:sens} runs it. The engagement-over-activation ranking --- the ranking
\cref{sec:med:euro} just recommended --- survives only while the hidden confounder explains less than
\pc{\nbFourFlipMed} of the fitted edge, 90\,\% interval $[\pc{\nbFourFlipLo},\ \pc{\nbFourFlipHi}]$,
with $\Prob(\text{no flip anywhere in the sweep}) = \nbFourPNoflip$. Half an edge is not an
outlandish amount of confounding for a behavioural mediator in product analytics. \textbf{If domain
knowledge cannot rule it out, the engagement-versus-activation ranking should be reported as
undecided} --- and no interval in \cref{tab:med:compare} or \cref{fig:med:forest} would ever have
told you so.

\input{build/floats/nb04_sensitivity}

Onboarding's story, by contrast, survives the same sweep. Its per-SD value falls to
\eur{\nbFourOnbValAtFlip} at the median flip point, down \pc{\nbFourOnbValDropPct} from
\eur{\nbFourValOnboarding} --- a real haircut, but not a reversal. The reason is structural:
onboarding's identification leans on the \emph{treatment} side (\cref{as:med:si1}), which is
exogenous here, whereas the engagement lever leans on the \emph{mediator} side
(\cref{as:med:si2}), which is not. \textbf{The fragile conclusions are exactly the ones that rest on
the untestable assumption}, and the sweep is what makes that visible.

\subsection*{Post-treatment variables: block a path or open one}

\Cref{sec:med:classical} showed conditioning on a mediator \emph{blocking} the path it was supposed
to measure. The complementary error is worse, because it moves the estimate in an unpredictable
direction: conditioning on a \emph{collider} --- a common effect of the treatment and something else
--- \emph{opens} a spurious path where none existed \citep{pearl2009, cole2002}. If the analyst
adjusts for ``opened the welcome email'', a variable caused by onboarding and also by an unlogged
engagement trait that separately drives conversion, then conditioning on it induces an association
between onboarding and that trait, and the estimate is biased by a path the DAG never contained.

The rule that covers both cases is the one to leave the room with: \textbf{a post-treatment variable
is never a free control.} Adjusting for it either removes the mechanism being priced or manufactures
a spurious one, and which of the two happened cannot be read off the regression output. Only the
graph knows, which is why the graph is drawn first.

\subsection*{Mediator measurement error, and the direct effect that is not there}

If ``engagement'' is a noisy proxy for the behaviour that actually drives conversion --- and in
product analytics it always is --- then the fitted $b$ is attenuated toward zero, the natural
indirect effect is \emph{understated}, and the share that has gone missing silently reappears as
\emph{direct} effect. A near-zero NDE like this chapter's is therefore doubly reassuring: attenuation
would have pushed it \emph{away} from zero, and it did not budge. The converse is the practical
warning: \textbf{a large estimated direct effect on real data should be interrogated for mediator
noise before it is interpreted as a mechanism the funnel model has missed.}

\subsection*{The limits of a linear decomposition}

Linearity is an \emph{economic} assumption, not merely a statistical one. A linear model prices the
tenth standard deviation of engagement exactly like the first --- no diminishing returns --- which is
the first thing a CMO will push back on, and rightly. Every euro-per-SD figure in
\cref{sec:med:euro} should be read as a \emph{local, first-SD} number. And if the return to
engagement depends on how well the user was onboarded (\cref{as:med:nointer} fails), then
$\text{NDE} + \text{NIE} \neq \text{TE}$, the clean two-way split of \cref{eq:med:te} stops adding
up, and \citeauthor{vanderweele2014}'s four-way decomposition is required
\citep{vanderweele2014, vanderweele2015}.

% =====================================================================================
\section{Takeaways}
\label{sec:med:take}

\begin{enumerate}[leftmargin=*, itemsep=.45em]
  \item \textbf{The total effect is not the decision; the split is.} Onboarding's effect on
        conversion is \nbFourBayesTotal{} and \pc{\nbFourPropMediated} of it is mediated
        (90\,\% $[\pc{\nbFourPropLo},\ \pc{\nbFourPropHi}]$), with a direct effect of
        \nbFourBayesDirect{} whose interval contains zero. It is a root investment whose payoff lands
        two stages downstream, and any attribution system that scores it on direct conversions prices
        it at \nbFourNaiveOnCoef{} and defunds it.

  \item \textbf{Never condition on a variable that lies on the path you are measuring.} Sliding
        activation into the outcome regression --- an ``obvious control'' --- cuts the
        engagement-to-conversion coefficient from \nbFourBlockingBOk{} to \nbFourBlockingBBad{} and
        would report an indirect effect of \nbFourBlockingNie{} instead of \nbFourBkNie. And the
        complementary error, conditioning on a collider, \emph{opens} a path rather than closing one.
        A post-treatment variable is never a free control.

  \item \textbf{The classical baseline is two regressions and a multiplication, and it is right.}
        \Cref{eq:med:product} is the mediation formula's linear special case, not a rival theory.
        \Citeauthor{baron1986}'s arithmetic under \citeauthor{imai2010}'s licence answers the whole
        question in milliseconds, and \nbFourClNCover{} of \nbFourNEstimands{} bootstrap intervals
        cover their planted value.

  \item \textbf{The bootstrap was worth running, and it barely mattered.} The indirect effect is a
        \emph{product}, whose sampling distribution is not normal in general --- but with
        $n = \nbFourNObs$ and both links far from zero ($t = \nbFourBkAT$, $t = \nbFourBkBT$), the
        product's bootstrap skew is \nbFourNieBootSkew{} and the delta method's linearisation
        (\cref{eq:med:delta}) lands within \pc{\nbFourSeGapPct} of it. \textbf{The delta method does
        not fail here.} It fails when a link is weak --- exactly the marginal case in which the
        analyst most needs an honest interval, and exactly why one bootstraps by default and is
        pleased, not disappointed, when it agrees.

  \item \textbf{The Bayesian arm got the interval for free, and the interval was not the problem.}
        The NIE is a deterministic function of two sampled coefficients (\cref{eq:med:derived}), so
        the posterior needs no delta method and no bootstrap --- multiply the draws. That is genuine
        convenience, and it is \emph{convenience, not capability}: the credible intervals came back
        slightly \emph{wider} (ratios \nbFourWidthMin--\nbFourWidthMax) and the point estimates
        differ from the classical ones by at most \nbFourMaxGap. The Bayesian machinery did not move
        the number and it did not sharpen the interval.

  \item \textbf{Randomizing the treatment does not identify the mechanism.} \Cref{as:med:si1} is what
        a coin flip buys. \Cref{as:med:si2} is not, because a coin flip on onboarding does not
        randomize engagement --- engagement is a behaviour, not a knob. A mediation analysis inside a
        perfect RCT is an observational study in an experiment's clothes.

  \item \textbf{In mediation, the dominant uncertainty is not in the interval.} Delete the mediators'
        common cause and both arms break \emph{identically} (\cref{tab:med:failfix}; they differ by
        \nbFourBadArmGap), each reporting a tight interval around a confounded number. A prior cannot
        see an unmeasured confounder, and neither can a bootstrap.

  \item \textbf{So price the assumption, since you cannot test it.} A hidden mediator--outcome
        confounder explaining \pc{\nbFourFlipMed} of the engagement-to-conversion edge
        (90\,\% $[\pc{\nbFourFlipLo},\ \pc{\nbFourFlipHi}]$) reverses the investment ranking that
        \cref{sec:med:euro} recommended --- while onboarding's near-fully-mediated story, which rests
        on the exogenous treatment side, survives with a \pc{\nbFourOnbValDropPct} haircut. The
        fragile conclusions are the ones resting on the untestable assumption, and the sweep is the
        only instrument that says which those are. The follow-up study is not a fancier model: it is
        an experiment that \emph{randomizes onboarding and logs the mediators}.
\end{enumerate}

\notebooknote{%
  This chapter is \code{notebooks/04\_funnel\_mediation.ipynb}. It runs in a fast teaching mode
  (\code{CMP\_FAST=1}, short chains) and a full mode (\code{CMP\_FAST=0}); every number here comes
  from the full run, emitted by \code{cmp.report} and injected as a macro. The simulator is
  \code{cmp.dgp.funnel} (seed \nbFourSeed), the classical arm is
  \code{cmp.classical.mediation\_product}, and the chain model is fitted with \code{pathmc}, which
  compiles the structural spec of \crefrange{eq:med:lik1}{eq:med:lik3} into PyMC and returns
  path-specific effects as posterior products of the coefficients along each route ---
  \cref{eq:med:product}, exactly, rather than a Monte-Carlo simulation through the graph.
  \Citet{imai2011} is the canonical treatment of what an experiment can and cannot buy here;
  \citet{vanderweele2015} is the reference to keep on the desk.}
